%%%PAGE 161 rot=0 legibility=good summary=磁场复合场中的摆线运动：E与B叠加求最大速度、配速法、摆线性质(h,T,S)、洛伦兹力分量冲量；圆约束小记；静止释放带电小球综合题
% 页面顶部边缘(y<0.05)露出的是上方另一页的文字残片，已忽略
%%%BLOCK id=161-01 ch=11 sec=摆线运动·最大速度 kind=derivation alt=none cont=none y=0.04-0.25
$V_x=ky$ 与 $E$ 无关，求 $V_{\max}$

%FIG p161_a 0.31 0.04 0.60 0.145
\notefig[0.3]{figures/p161_a.png}
\[ qEy_m=\tfrac12mV_m^2-\tfrac12mV_0^2 \]
$V_m=ky_m$\quad \unclear{令}$E=0$\quad 又$V_0=kR_0$\quad 故 \struck{$V_m=kR_0$}\quad $y_0=R_0$

% 下一行 k=(q/m)B 之后有涂改液痕迹
\[ k=\frac{q}{m}B\qquad V_m^2-\frac{2E}{B}V_m-V_0^2=0\qquad V_m=\frac{E}{B}\pm\sqrt{\left(\frac{E}{B}\right)^2+V_0^2} \]
%%%END
%%%BLOCK id=161-02 ch=11 sec=摆线运动·配速法 kind=method alt=none cont=none y=0.25-0.59
\textbf{〈配速法〉}

%FIG p161_b 0.13 0.27 0.40 0.405
\notefig[0.3]{figures/p161_b.png}
\[ V_0=\frac{mg}{qB}\qquad 0=1+(-1). \]
\[ \text{两个分运动}\left\{\begin{array}{l}\text{向右匀速直线运动}\\ \text{逆时针圆周运动}\end{array}\right.\ \to\ \text{摆线} \]

\textbf{轨迹为摆线}
\[ V_{\max}=V_0+V_0=\frac{2mg}{qB}\qquad V_{\min}=V_0-V_0=0 \]
竖直最大位移\quad$mgh=\tfrac12mV_{\max}^2$\quad$h=\dfrac{2m^2g}{q^2B^2}$\quad$R=\dfrac{mV_0}{qB}=\dfrac{m^2g}{q^2B^2}$\quad$\therefore h=2R$ 为直径长度

$T=\dfrac{2\pi m}{qB}$\quad 到最低点时间\quad$t=\dfrac{T}{2}=\dfrac{\pi m}{qB}$

1个周期内前进距离\quad$S=V_0T=\dfrac{2\pi m^2g}{q^2B^2}$
%%%END
%%%BLOCK id=161-03 ch=11 sec=摆线运动·洛伦兹力分量冲量 kind=derivation alt=none cont=none y=0.575-0.70
%FIG p161_c 0.10 0.575 0.285 0.655
\notefig[0.25]{figures/p161_c.png}
\[
mgt\ \left\{\begin{aligned}
&mgh=\tfrac12mV_2^2-\tfrac12mV_1^2\\
&qB\Delta x=m\Delta v_y\\
&qB\Delta y\,\text{\struck{$+mgt$}}=m\Delta v_x
\end{aligned}\right.
\]
〈洛伦兹分力冲\unclear{量}〉 % 末字上有涂改
%%%END
%%%BLOCK id=161-04 ch=11 sec=有界磁场·圆约束 kind=method alt=none cont=none y=0.71-0.80
% 原稿此块被大圈（云形曲线）圈起，左上角“圆约束”三字另有方框
\begin{keybox}[圆约束]
两次进磁区圆心\unclear{角}$150^\circ$，经$m$次偏转后再入射
\[ m\cdot150^\circ=n\cdot360^\circ \]
\[ m=12\qquad n=5\qquad \text{求最小正整数} \]
\end{keybox}
%%%END
%%%BLOCK id=161-05 ch=11 sec=摆线运动·综合题 kind=problem alt=06,09 cont=none y=0.77-1.00
\begin{problem}
%FIG p161_d 0.01 0.76 0.30 0.955
\notefig[0.3]{figures/p161_d.png}
静止释放，在最低点曲率半径为$2y_m$ % 此行“静止释放”处有涂改液痕迹

① 求小球运动$P(x,y)$速度
\begin{solution}
\[ mgy=\tfrac12mv^2\qquad v=\sqrt{2gh} \]
\end{solution}
② 求第一次下降的最大距离$y_m$
\begin{solution}
\[ qvB-mg=\frac{mv^2}{R}\qquad v_1=\sqrt{2gy_m} \]
\[ R=2y_m\qquad \therefore y_m=\frac{2m^2g}{q^2B^2} \]
\end{solution}
③ 在竖直方向加一匀强电场$E$（$E>\dfrac{mg}{q}$），求$V_{\max}$ % 原稿“匀强”二字插写在“一”字上方
\begin{solution}
\[ (qE-mg)y_m=\tfrac12mV_m^2 \]
\[ qV_mB\ \unclear{\pm}\ mg+qE=\frac{mV_m^2}{R}\qquad R=2y_m \] % CHECK: 此式中间的符号有涂改（像 + 与 - 叠写）；按解得结果应为 qV_mB+mg-qE=mV_m^2/R
解得 $V_m=\dfrac{2}{qB}(qE-mg)$
\end{solution}
\end{problem}
%%%END
%%%PAGEEND 161
